% The four parts of the Day 4 lab. The steps follow Labs/day04/README.md.
%
% Hardware status: not tested on hardware. A value that only a board can
% give is written as "..." and the students measure it. The deck does not
% show a result that the students must predict.

% ------------------------------------------------------------------- Part A
\begin{frame}{Part A: quantize by hand (30 min)}
  \small
  Run sections 0 and 1 of the notebook. Then write three functions in
  section 2.
  \vskip 0.3em
  \renewcommand{\arraystretch}{1.25}
  \begin{tabular}{@{}L{0.09\textwidth}L{0.36\textwidth}L{0.45\textwidth}@{}}
    \toprule
    \textbf{Task} & \textbf{You write} & \textbf{The notebook checks} \\
    \midrule
    1 & \code{scale\_zero\_point}: the scale and the zero point of a range
      & The range $[-1.00, 1.55]$ gives $S = 0.01$ and $Z = -28$ \\
    2 & \code{quantize} and \code{dequantize}: one line for each function
      & The five numbers of the lecture exercise give $-28$, 22, $-61$,
        95, and 127 \\
    3 & Your prediction: the range with the smallest error
      & Not checked. You compare it with the result. \\
    \bottomrule
  \end{tabular}
  \vskip 0.5em
  {\footnotesize
  \[
    S = \frac{r_{\max} - r_{\min}}{255} \quad\;
    Z = \text{round}\!\left(-128 - \frac{r_{\min}}{S}\right) \quad\;
    q = \text{clamp}\!\left(\text{round}\!\left(\frac{r}{S}\right) + Z\right)
  \]}
  \begin{itemize}
    \item Use \code{np.round} and \code{np.clip}. The clamp limits the
          result to the range from $-128$ to 127. The way back is
          $\hat{r} = S\,(q - Z)$.
  \end{itemize}
\end{frame}

\begin{frame}{Part A: the rounding error and the clipping error}
  \small
  \begin{columns}[T]
    \begin{column}{0.47\textwidth}
      \begin{block}{Section 3: the rounding error}
        \begin{itemize}
          \item The code quantizes all acceleration values of your
                training windows.
          \item \textbf{You see:} the largest error and the typical error.
          \item Check the two rules of the lecture: $S/2$ and
                $S/\sqrt{12}$.
          \item Then the same range with 7, 6, 4, and 2 bits.
        \end{itemize}
      \end{block}
    \end{column}
    \begin{column}{0.49\textwidth}
      \begin{block}{Section 4: the clipping error}
        \begin{itemize}
          \item Four ranges for the same tensor: the minimum and the
                maximum, and the ranges that hold 99.9, 99, and 90 percent
                of the values.
          \item \textbf{You see:} the scale, the clipped values, and the
                RMS error, with 8 bits and with 4 bits.
        \end{itemize}
      \end{block}
    \end{column}
  \end{columns}
  \vskip 0.3em
  \begin{itemize}
    \item \textbf{You record:} the scale, the zero point, the two errors,
          and the answers to the two questions of section 4.
  \end{itemize}
  \begin{exerciseblock}
    \textbf{Prediction (task 3).} Which of the four ranges gives the
    smallest RMS error with 8 bits? Which range with 4 bits? Write the
    prediction before you run the experiment.
  \end{exerciseblock}
\end{frame}

% ------------------------------------------------------------------- Part B
\begin{frame}{Part B: post-training quantization (50 min)}
  \begingroup\centering
  \begin{tikzpicture}[font=\footnotesize, >={Stealth[length=5pt]},
      b/.style={draw=coolgray, rounded corners=3pt, fill=boxgray,
                minimum width=1.75cm, minimum height=1.0cm, align=center},
      lab/.style={font=\scriptsize, text=coolgray, align=center}]
    \node[b] (a) at (0,0)    {Float\\CNN};
    \node[b] (b) at (2.15,0) {Calibration\\set};
    \node[b, draw=cardinalred] (c) at (4.3,0)  {Convert};
    \node[b, draw=cardinalred] (d) at (6.45,0) {Test};
    \node[b] (e) at (8.6,0)  {Look\\inside};
    \foreach \a/\b in {a/b, b/c, c/d, d/e} { \draw[->] (\a) -- (\b); }
    \node[lab] at (0,-0.95)    {section 5\\852 parameters};
    \node[lab] at (2.15,-0.95) {50 windows of\\each class};
    \node[lab] at (4.3,-0.95)  {section 6\\task 4};
    \node[lab] at (6.45,-0.95) {section 7\\task 5};
    \node[lab] at (8.6,-0.95)  {section 8};
  \end{tikzpicture}\par\endgroup
  \vskip 0.5em
  \small
  \renewcommand{\arraystretch}{1.25}
  \begin{tabular}{@{}L{0.09\textwidth}L{0.40\textwidth}L{0.41\textwidth}@{}}
    \toprule
    \textbf{Task} & \textbf{You write} & \textbf{The notebook checks} \\
    \midrule
    4 & \code{convert\_int8}: the generator of the calibration inputs, and
        the five settings of the converter
      & The input tensor of the model has the type \code{int8} \\
    5 & \code{quantize\_input}: the \code{int8} input of the model, and the
        number of clipped values
      & The numbers of the lecture exercise, with two clipped values \\
    \bottomrule
  \end{tabular}
  \vskip 0.4em
  \begin{itemize}
    \item The generator gives one input for each \code{yield}: a list with
          one array of the shape \code{(1, 100, 3)}.
  \end{itemize}
\end{frame}

\begin{frame}[fragile]{Part B: test the model, and look inside}
  \small
  Section 7 prints one row for each model. The accuracy and the clipped
  input values are in percent:
\begin{lstlisting}
model                   bytes  accuracy    idle terrest    lift maritim  clipped
float32                   ...       ...     ...     ...     ...     ...     0.00
int8, all classes         ...       ...     ...     ...     ...     ...      ...
\end{lstlisting}
  \begin{itemize}
    \item \textbf{You record:} the two rows, the largest difference between
          the two outputs, and the next step by the rule of the lecture.
    \item The file does not become four times smaller. Find the reason in
          the output of section 6.
  \end{itemize}
  \vskip 0.2em
  Section 8 prints each tensor of the \code{int8} model that has a scale.
  Answer three questions in the report:
  \begin{enumerate}
    \item Which scale and which zero point has the input tensor? Compare
          them with the result of your function for the range of the
          calibration set.
    \item How many scales has the weight tensor of each convolution, and
          which zero point?
    \item Which type have the biases?
  \end{enumerate}
\end{frame}

\begin{frame}{Part B: change the calibration set}
  \small
  The conversion gives no error when the calibration set is wrong. Section
  9 makes this error on purpose.
  \vskip 0.3em
  \renewcommand{\arraystretch}{1.25}
  \begin{tabular}{@{}L{0.20\textwidth}L{0.36\textwidth}L{0.34\textwidth}@{}}
    \toprule
    \textbf{Experiment} & \textbf{Calibration set} & \textbf{You record} \\
    \midrule
    1 & 50 windows of one class only. One conversion for each class.
      & The accuracy, the accuracy of each class, and the clipped input
        values \\
    2 & All classes, with 200, 20, and 4 windows
      & The accuracy and the clipped input values \\
    \bottomrule
  \end{tabular}
  \vskip 0.5em
  \begin{itemize}
    \item The table of section 1 gives the smallest value and the largest
          value of each class. Use it for your prediction.
    \item The column \code{clipped} needs no labels. It shows the problem
          also for inputs with no known class.
  \end{itemize}
  \begin{exerciseblock}
    \textbf{Prediction (task 6).} Which calibration class gives the lowest
    accuracy? Is there a class that gives no loss? Which is the smallest
    set of all classes with no loss: 4, 20, or 200 windows?
  \end{exerciseblock}
\end{frame}

% ------------------------------------------------------------------- Part C
\begin{frame}{Part C: quantization-aware training (40 min)}
  \begingroup\centering
  \begin{tikzpicture}[font=\footnotesize, >={Stealth[length=5pt]},
      b/.style={draw=coolgray, rounded corners=3pt, fill=boxgray,
                minimum width=1.75cm, minimum height=1.0cm, align=center},
      o/.style={b, draw=darkcyan, minimum width=1.3cm},
      lab/.style={font=\scriptsize, text=coolgray, align=center}]
    \node[o] (a) at (0,0)     {Observer};
    \node[b, draw=cardinalred] (b) at (1.95,0) {QuantConv1D};
    \node[o] (c) at (3.9,0)   {Observer};
    \node[lab] (d) at (5.2,0) {\dots};
    \node[b, draw=cardinalred] (e) at (6.6,0) {QuantDense};
    \node[o] (f) at (8.55,0)  {Observer};
    \foreach \a/\b in {a/b, b/c, c/d, d/e, e/f} { \draw[->] (\a) -- (\b); }
    \node[lab] at (0,-0.95)    {the input\\8 bits};
    \node[lab] at (1.95,-0.95) {weights with\\8, 4, 3, or 2 bits};
    \node[lab] at (3.9,-0.95)  {an activation\\8 bits};
    \node[lab] at (6.6,-0.95)  {weights with\\8, 4, 3, or 2 bits};
    \node[lab] at (8.55,-0.95) {before the\\softmax};
  \end{tikzpicture}\par\endgroup
  \vskip 0.4em
  \small
  \begin{itemize}
    \item The notebook builds the CNN again with layers that quantize and
          dequantize in the forward pass. All layers use your function
          \code{fake\_quant}.
    \item \textbf{Task 7.} Replace the last line of \code{fake\_quant} with
          the straight-through estimator. Use \code{ops.stop\_gradient}.
    \item The notebook checks two facts: the forward value is the quantized
          value, and the gradient is 1 for each value.
    \item \textbf{You record:} the accuracy with simulated quantization
          with 8 bits. It must be near the accuracy of the \code{int8}
          LiteRT model of Part B.
  \end{itemize}
  \keypoint{With no straight-through estimator, the gradient is
    \code{None}, and the training cannot start.}
\end{frame}

\begin{frame}[fragile]{Part C: fewer bits, the layer, and the real model}
  \small
  Section 11 trains the model four times. The activations keep 8 bits:
\begin{lstlisting}
weight bits  values  no new training  quantization-aware training
          8     255             ...%                         ...%
          4      15             ...%                         ...%
          3       7             ...%                         ...%
          2       3             ...%                         ...%
\end{lstlisting}
  \begin{itemize}
    \item ``No new training'' is post-training quantization. LiteRT has no
          format below 8 bits: the rows for 4, 3, and 2 bits are a
          simulation only.
    \item \textbf{Section 12:} one layer has weights with 4 bits or with 2
          bits. All other values stay in float. \textbf{You record:} the
          loss of each layer, and the layer that you keep at 8 bits.
    \item \textbf{Section 13:} the model that trained with 8 bits becomes a
          real \code{int8} file. \textbf{You record:} the three rows of
          the table.
  \end{itemize}
  \begin{exerciseblock}
    \textbf{Prediction (task 8).} Which accuracy has the model with weights
    of 2 bits and no new training? Does quantization-aware training repair
    the loss?
  \end{exerciseblock}
\end{frame}

% ------------------------------------------------------------------- Part D
\begin{frame}{Part D: on the board (30 min)}
  \small
  \renewcommand{\arraystretch}{1.2}
  \begin{tabular}{@{}L{0.07\textwidth}L{0.70\textwidth}L{0.12\textwidth}@{}}
    \toprule
    \textbf{Step} & \textbf{Work} & \textbf{Time} \\
    \midrule
    1 & Run sections 14 and 15 of the notebook. They train the feature
        model of Day 3 and convert it two times. & 5 min \\
    2 & Write your prediction in the report. & 2 min \\
    3 & Upload the sketch with the float model. Record the numbers.
      & 6 min \\
    4 & Complete tasks D1 and D2. Upload the sketch with the \code{int8}
        model. Record the numbers. & 8 min \\
    5 & Make the four motions. After the test set, the sketch classifies
        the motion of the kit, as on Day 3. & 4 min \\
    6 & Complete the comparison table. Write the Decision Log. & 5 min \\
    \bottomrule
  \end{tabular}
  \vskip 0.5em
  \begin{itemize}
    \item The notebook writes four files into
          \code{sketches/motion\_quant/}: \code{model\_float.h},
          \code{model\_int8.h}, \code{model\_settings.h}, and
          \code{test\_set.h}.
    \item \code{test\_set.h} holds 80 test windows, with the class that
          each model gives on the laptop.
  \end{itemize}
\end{frame}

\begin{frame}{Part D: tasks D1 and D2}
  \small
  Open the sketch \code{motion\_quant.ino}. Set
  \code{Tools > PSRAM > Disabled}. The line \code{\#define MODEL\_INT8 0}
  selects the float model. With the value 1, it selects the \code{int8}
  model.
  \vskip 0.3em
  \renewcommand{\arraystretch}{1.25}
  \begin{tabular}{@{}L{0.09\textwidth}L{0.84\textwidth}@{}}
    \toprule
    \textbf{Task} & \textbf{You write, in the function \code{classify()}} \\
    \midrule
    D1 & Quantize the input with the scale and the zero point of
         \code{input->params}. Calculate
         \code{lroundf(x / scale) + zero\_point}. Limit the result to the
         range from $-128$ to 127. Write it into
         \code{input->data.int8[i]}. Count the clipped values. \\
    D2 & Dequantize the output with the scale and the zero point of
         \code{result->params}. Write
         \code{scale * (result->data.int8[k] - zero\_point)} into
         \code{output[k]}. \\
    \bottomrule
  \end{tabular}
  \vskip 0.5em
  \begin{itemize}
    \item The float model needs no change. Upload it first.
    \item When task D1 is complete, remove the line
          \code{task\_d1\_complete = false;}. Do the same for task D2.
    \item Before that, the sketch with the \code{int8} model prints that
          the tasks are not complete, and it stops.
  \end{itemize}
\end{frame}

\begin{frame}[fragile]{Part D: the output of the sketch}
  \small
  The Serial Monitor at 115200 baud shows these lines for the \code{int8}
  model:
\begin{lstlisting}
Model:                int8
Model size (flash):   ... bytes
Sketch size (flash):  ... bytes
Arena size:           4096 bytes
Arena used:           ... bytes
Input tensor:  type INT8, 63 bytes, scale ..., zero point ...
Output tensor: type INT8, 4 bytes, scale 0.003906, zero point -128
Test set:             80 windows
Test set, correct:    ... of 80 (... percent)
Test set, same class as on the laptop: ... of 80
Test set, clipped input values: ... of 5040 (laptop: ...)
Test set, invoke_us:  median ..., largest ...
\end{lstlisting}
  \begin{itemize}
    \item The sketch runs the 80 test windows on the board. The line
          \code{same class as on the laptop} must show 80 of 80. If it
          does not, check tasks D1 and D2.
    \item The float model prints the same lines, with the type
          \code{FLOAT32} and with no line for the clipped values.
  \end{itemize}
  \source{The output shows the form of the lines. Nobody measured the
    values on the kit. The output scale is the scale of a softmax in
    \texttt{int8}: 1/256.}
\end{frame}

\begin{frame}{Part D: the comparison table}
  \small
  \renewcommand{\arraystretch}{1.25}
  \begin{tabular}{@{}L{0.30\textwidth}L{0.62\textwidth}@{}}
    \toprule
    \textbf{Number} & \textbf{Where you read it, for each model} \\
    \midrule
    File size, test accuracy, peak of the activations
                    & Section 14 of the notebook \\
    Flash           & \code{Sketch uses ... bytes} of the build \\
    Arena           & \code{Arena used} of the Serial Monitor \\
    Accuracy on the board & \code{Test set, correct} \\
    Agreement with the laptop & \code{Test set, same class as on the
                      laptop} \\
    Latency         & \code{Test set, invoke\_us}: the median and the
                      largest value \\
    \bottomrule
  \end{tabular}
  \vskip 0.5em
  \vskip 0.3em
  Write the two values and the change of each number. Then compare each
  result with your prediction.
  \begin{exerciseblock}
    \textbf{Prediction.} How does the \code{int8} model change the flash
    use, the arena, and the time of \code{Invoke()}? The ESP32-S3 has a
    circuit for float arithmetic, and the library has no optimized kernels
    for this chip.
  \end{exerciseblock}
\end{frame}
